% page 452 of the facsimile (image p-451 of the scan)
% block 160.0 x 242.0 mm ; left margin 8.0 mm
\pgc{-12.700mm}{0.000mm}{
\l{\cel{3}444}
\l{}
\l{planko. Ligno-peco, fendita longesale, plu longa kam larja,}
\l{\cel{3}e nur poke dina. - DEFIS.}
\l{}
\l{\sou{\textquotedbl{}plankton\textquotedbl{}.}\cel{2}I. Ensemblo ek omna vivanti (animali e veje-}
\l{\cel{3}tanti) qui fluktuas en la aquo di maro o di lago, inter}
\l{\cel{3}la surfaco e la fundo. - II. Animali o vejetanti tre mi-}
\l{\cel{3}kra di la aquo di maro o di lago.- DEFIRS.}
\l{}
\l{\sou{planto}. (\sou{bot.}) Omno quo fixigesas en sulo per radiki.- DEFIS}
\l{}
\l{\sou{plantaca}r. (\sou{trans.}) Fixigar en sulo per radiki (vejetanto).}
\l{\cel{3}- DEFIS.}
\l{}
\l{\sou{plantago}. (\sou{bot.}) Planto herbatra di qua la spiko kovresas ye}
\l{\cel{3}mikra semini - tipo di la familio \textquotedbl{}plantaginei\textquotedbl{}. -}
\l{\cel{3}L.\cel{1}\sou{plantago}. - EFI.}
\l{}
\l{\sou{plantigrado}. (\sou{zool.}) Mamifero qua, marchante, pozas la plando}
\l{\cel{3}di sua pedi sur sulo. (Ex.: urso, daxo, e c.) - DEFIS.}
\l{}
\l{\sou{plantono}.\cel{2}(\sou{armeo}).Soldato qua stacas avan la pordo di oficiro}
\l{\cel{3}alta-grada, por transmisar komandi, e c. - F.}
\l{}
\l{\sou{plaso}. I. Spaco, okupata od okupebla, da singla persono en}
\l{\cel{3}irga loko (veturo, teatro, cinemo, e c.) - II. Loko propra,}
\l{\cel{3}rezervita o konvenanta por singla kozo.- III. Rolo atri-}
\l{\cel{3}buita ad employato en entraprezuro, komerco-kontori, fa-}
\l{\cel{3}brikeyo, e c. - DEF.}
\l{}
\l{\sou{plasmo}. (\sou{biol.}) Parto liquida di sango, en qua fluktuas la}
\l{\cel{3}globuli di sango. - DEFIS.}
\l{}
\l{plasmodermo. (\sou{biol.}) Materio protoplasma qua formacas la}
\l{\cel{3}komuniki intercelula quin onu vidas che la vejetanti e}
\l{\cel{3}che ula animali. - DEFIRS.}
\l{}
\l{\sou{plasmodio}. (\sou{parazitologio}). Parazito di la genero \textquotedbl{}plasmo-}
\l{\cel{3}dium\textquotedbl{}, en qua asemblesas la tri speci de hematozoi di}
\l{\cel{3}la marcho-febro amiboida, e qua vivas en la hematii di la}
\l{\cel{3}homo e di la animali, e genitesas per skizogonio o sporo-}
\l{\cel{3}gonio. - DEFIRS.}
\l{}
\l{\sou{plasmolizo.}\cel{2}(\sou{biol}.)Fenomeno osmosala, qua abutas a la destruk-}
\l{\cel{3}to di la plasmo. - DEFIRS.}
\l{}
\l{\sou{plasmosomo.}\cel{2}(\sou{histologio}).Nukleolo vera, nedependanta de la}
\l{\cel{3}reto de linino, qua povas movar su sur la nukleo e kolor-}
\l{\cel{3}izesas intense da la kolorizivo di plasmo. - DEFIRS.}
\l{}
\l{plastido..(\sou{biol}.)Ensemblo ek la substanci vivanta di la celulo,}
\l{\cel{3}ecepte la rezervo-substanci, la ek-sekrecaji, la envelo-}
\l{\cel{3}po-membrano. - DEFIS.}
\l{}
\l{plastica. Quan onu povas formizar o modlar.}
}
